\subsection{\bilingual{Proof of the Polynomial Decay of the Left Tail}{左尾部多项式衰减的证明}}

\bilingual{
    We restate the expression of \eqref{eq:4.16} for the case of $a\in[0,1]$ as follows: for any $t>0$, $\zeta<-a+\frac{\pi^2}2\sigma^2t$, there exists some $\varepsilon=\varepsilon(a,t,\zeta)>0$, some $N_0=N_0(a,t,\zeta)>0$, such that for all $N\ge N_0$,
}{
    现将 \eqref{eq:4.16} 重述: 对任意 $t>0$, $\zeta<-a+\frac{\pi^2}2\sigma^2t$, 存在某个 $\varepsilon=\varepsilon(a,t,\zeta)>0$, 某个 $N_0=N_0(a,t,\zeta)>0$, 使得对于所有 $N\ge N_0$,
}
\begin{equation}
    \mathbb P\left(
        \mathcal M^{N|a}_{\ell_N}
        \le
        \zeta\log N
    \right)
    \le
    N^{-\varepsilon}.
    \tag{\ref{eq:4.16}}
\end{equation}

\begin{proof}
    \quad

    \bilingual{
        Now we proceed to show \eqref{eq:4.16}. Let $b>1$ in this part. Let
    }{
        现在继续证明 \eqref{eq:4.16}. 在本部分中设 $b>1$. 记
    }
    \begin{equation*}
        F_N
        :=
        \left\{
            \exists j,n\in[\![0,\ell_N]\!],\,
            \exists |u|=n:\,
            j\le n,\,
            \mathcal V^{N|a}(u_j)-\mathcal V^{N|a}(u)
            >
            b\log N
        \right\}
    \end{equation*}
    \bilingual{
        be the event that there exists a position drop of over $b\log N$ from an ancestor to its descendant. This event can also be written as
    }{
        表示存在从祖先到其后代的位移下降 (历史回撤) 超过 $b\log N$ 的事件. 该事件也可写为
    }
    \begin{equation*}
        F_N
        =
        \left\{
            \begin{aligned}
                \exists j\in[\![0,\ell_N-1]\!],\,
                \exists k\in[\![1,\#\mathcal Z^{N|a}_j]\!],
                \\ \exists n\in[\![j+1,\ell_N]\!],\,
                \exists |u|=n:
                \\
                    u_j=\nu^{(k)}_j,
                \\
                    \mathcal V^{N|a}(u_j)
                    -
                    \mathcal V^{N|a}(u_i)
                    \le
                    b\log N,\,
                    \forall i\in[\![j+1,n-1]\!],
                \\
                    \mathcal V^{N|a}(u_j)
                    -
                    \mathcal V^{N|a}(u)
                    >
                    b\log N
            \end{aligned}
        \right\},
    \end{equation*}
    \bilingual{
        where we denote
    }{
        其中记
    }
    \begin{equation*}
        \mathcal Z^{N|a}_j
        =:
        \left\{
            \nu^{(k)}_j:\,
            k\in\left[\!\!\left[1,\#\mathcal Z^{N|a}_j\right]\!\!\right]
        \right\}.
    \end{equation*}
    \bilingual{
        Calculate first the upper bound of the probability of this event. By the union bound for $j$, \nameref{lemma:branching-property} for $\mathcal V$, and \nameref{lemma:many-to-one} for $\mathcal V$. For each fixed $j$, the first-crossing events at distinct descendant generations are disjoint. After applying many-to-one, write $A_{j,n}$ for the corresponding spine event at elapsed time $n$.
    }{
        先计算这个事件的概率上界. 通过先对 $j$ 使用并集界, 再对 $\mathcal V$ 使用\nameref{lemma:branching-property}, 然后对 $\mathcal V$ 使用\nameref{lemma:many-to-one}. 对每个固定的 $j$, 不同后代代数处的首次越界事件两两不交. 应用多对一公式后, 记 $A_{j,n}$ 为经过 $n$ 步的对应脊柱事件.
    }
    \begin{align*}
        \mathbb P(F_N)
        &
        \le
        \sum_{j=0}^{\ell_N-1}
        \mathbb E\left[
            \mathbb E\left(
                \sum_{k=1}^{\#\mathcal Z^{N|a}_j}
                {\bf 1}{
                    \left\{
                        \text{\fontsize{9}{0} $
                        \begin{aligned}
                            \exists n\in[\![j+1,\ell_N]\!],\,
                            \exists |u|=n:\,
                            u_j=\nu^{(k)}_j,
                            \\[-8pt]
                                \mathcal V^{N|a}(u_j)
                                -
                                \mathcal V^{N|a}(u_i)
                                \le
                                b\log N,\\[-8pt] \,
                                \forall i\in[\![j+1,n-1]\!],
                            \\[-8pt]
                                \mathcal V^{N|a}(u_j)
                                -
                                \mathcal V^{N|a}(u)
                                >
                                b\log N
                        \end{aligned}
                        $}
                    \right\}
                }
                \middle|
                \mathscr F_j
            \right)
        \right]
        \\ &
        =
        \sum_{j=0}^{\ell_N-1}
        \!\!
        \mathbb E\!\left[
            \!
            \mathbb E\!\left(
                \sum_{m=1}^N
                \!\!
                {\bf 1}_{
                    \left\{
                        \#\mathcal Z^{N|a}_j
                        =
                        m
                    \right\}
                }
                \sum_{k=1}^m
                \!
                {\bf 1}{
                    \left\{
                        \text{\fontsize{9}{0} $
                        \begin{aligned}
                            \exists n\in[\![j+1,\ell_N]\!],\,
                            \exists |u|=n:\,
                            u_j=\nu^{(k)}_j,
                            \\[-8pt]
                                \mathcal V^{N|a}(u_j)
                                -
                                \mathcal V^{N|a}(u_i)
                                \le
                                b\log N,\\[-8pt]
                                \forall i\in[\![j+1,n-1]\!],
                            \\[-8pt]
                                \mathcal V^{N|a}(u_j)
                                -
                                \mathcal V^{N|a}(u)
                                >
                                b\log N
                        \end{aligned}
                        $}
                    \right\}
                }
                \middle|
                \mathscr F_j
            \right)
        \right]
        \\ &
        =
        \sum_{j=0}^{\ell_N-1}
        \!\!
        \mathbb E\!\left[
            \!
            \sum_{m=1}^N
            \!\!
            {\bf 1}_{
                \left\{
                    \#\mathcal Z^{N|a}_j
                    =
                    m
                \right\}
            }
            \sum_{k=1}^m
            \!
            \mathbb E\!\left(
                {\bf 1}{
                    \left\{
                        \text{\fontsize{9}{0} $
                        \begin{aligned}
                            \exists n\in[\![j+1,\ell_N]\!],\,
                            \exists |u|=n:\,
                            u_j=\nu^{(k)}_j,
                            \\[-8pt]
                                \mathcal V^{N|a}(u_j)
                                -
                                \mathcal V^{N|a}(u_i)
                                \le
                                b\log N,\\[-8pt]
                                \forall i\in[\![j+1,n-1]\!],
                            \\[-8pt]
                                \mathcal V^{N|a}(u_j)
                                -
                                \mathcal V^{N|a}(u)
                                >
                                b\log N
                        \end{aligned}
                        $}
                    \right\}
                }
                \middle|
                \mathscr F_j
            \right)
        \right]
        \\ &
        \le
        \sum_{j=0}^{\ell_N-1}
        \!\!
        \mathbb E\!\left[
            \!
            \sum_{m=1}^N
            \!\!
            {\bf 1}_{
                \left\{
                    \#\mathcal Z^{N|a}_j
                    =
                    m
                \right\}
            }
            \sum_{k=1}^m
            \!
            \mathbb E\!\left(
                {\bf 1}{
                    \left\{
                        \text{\fontsize{9}{0} $
                        \begin{aligned}
                            \exists n\in[\![j+1,\ell_N]\!],\,
                            \exists |u|=n:\,
                            u_j=\nu^{(k)}_j,
                            \\[-8pt]
                                \mathcal V(u_j)
                                -
                                \mathcal V(u_i)
                                \le
                                b\log N,\\[-8pt]
                                \forall i\in[\![j+1,n-1]\!],
                            \\[-8pt]
                                \mathcal V(u_j)
                                -
                                \mathcal V(u)
                                >
                                b\log N
                        \end{aligned}
                        $}
                    \right\}
                }
                \middle|
                \mathscr F_j
            \right)
        \right]
        \\ &
        =
        \sum_{j=0}^{\ell_N-1}
        \!\!
        \mathbb E\!\left[
            \!
            \sum_{m=1}^N
            \!\!
            {\bf 1}_{
                \left\{
                    \#\mathcal Z^{N|a}_j
                    =
                    m
                \right\}
            }
            \sum_{k=1}^m
            \!
            \mathbb P\left(
                \text{\fontsize{9}{0} $
                \begin{aligned}
                    \exists n\in[\![1,\ell_N-j]\!],\,
                    \exists |u|=n:
                    \\[-8pt]
                        \mathcal V(u_i)
                        \ge
                        -b\log N,\\[-8pt]
                        \forall i\in[\![1,n-1]\!],
                    \\[-8pt]
                        \mathcal V(u)
                        <
                        -b\log N
                \end{aligned}
                $}
            \right)
        \right]
        \\ &
        \le
        N
        \sum_{j=0}^{\ell_N-1}
        \mathbb P\left(
            \text{\fontsize{9}{0} $
            \begin{aligned}
                \exists n\in[\![1,\ell_N-j]\!],\,
                \exists |u|=n:
                \\[-8pt]
                    \mathcal V(u_i)
                    \ge
                    -b\log N,\\[-8pt]
                    \forall i\in[\![1,n-1]\!],
                \\[-8pt]
                    \mathcal V(u)
                    <
                    -b\log N
            \end{aligned}
            $}
        \right)
        \\ &
        \le
        N
        \sum_{j=0}^{\ell_N-1}
        \sum_{n=1}^{\ell_N-j}
        \mathbb E\left[
            \sum_{|u|=n}
            {\bf 1}_{
                \left\{
                    \mathcal V(u_i)
                    \ge
                    -b\log N,\,
                    \forall i\in[\![1,n-1]\!],\,
                    \mathcal V(u)
                    <
                    -b\log N
                \right\}
            }
        \right]
        \\ &
        =
        N
        \sum_{j=0}^{\ell_N-1}
        \sum_{n=1}^{\ell_N-j}
        \mathbb E\left[
            e^{Y_n}
            {\bf 1}_{
                \left\{
                    Y_i
                    \ge
                    -b\log N,\,
                    \forall i\in[\![1,n-1]\!],\,
                    Y_n
                    <
                    -b\log N
                \right\}
            }
        \right]
        \\ &
        \le
        N^{1-b}
        \sum_{j=0}^{\ell_N-1}
        \sum_{n=1}^{\ell_N-j}
        \mathbb P(A_{j,n})
        \\
        &=
        N^{1-b}
        \sum_{j=0}^{\ell_N-1}
        \mathbb P\left(\bigcup_{n=1}^{\ell_N-j} A_{j,n}\right)
        \\
        &\le N^{1-b}\ell_N
        \le N^{1-b}(t+o(1))\log^3N
        = N^{-(b-1)+o_N(1)}.
    \end{align*}
    \bilingual{
        First, the assumptions imply $\sigma^2>0$. The one-step spine law is
        the probability measure
        \[
            Q(A)=\mathbb E\left[\sum_i e^{-\Xi_i}{\bf1}_{\{\Xi_i\in A\}}\right].
        \]
        If its variance were zero, its centered law would be concentrated at
        zero. The many-to-one identity then gives
        \[
            \mathbb E\left[\sum_i e^{-\Xi_i}{\bf1}_{\{\Xi_i\ne0\}}\right]=0,
        \]
        so almost surely every child is at zero. Thus
        \[
            \mathbb E[\#\Xi]=\mathbb E\left[\sum_i e^{-\Xi_i}\right]=1,
        \]
        contradicting $\psi(0)>0$.

        Fix $a\in[0,1]$, $t>0$, and a target
        $\zeta<\mu$, where $\mu=-a+\frac{\pi^2}2\sigma^2t>-a$.
        At $b=1$ and $d=\frac{\pi^2}{2}\sigma^2$, the expression
        $a+\zeta-dtb^{-2}$ equals $\zeta-\mu<0$. By continuity, we may
        therefore choose $b>1$ close enough to $1$ and
        $d<\frac{\pi^2}{2}\sigma^2$ close enough to
        $\frac{\pi^2}{2}\sigma^2$ so that
        $a+\zeta-dtb^{-2}<0$ still holds.
    }{
        首先, 假设蕴含 $\sigma^2>0$. 一步脊柱律是如下概率测度:
        \[
            Q(A)=\mathbb E\left[\sum_i e^{-\Xi_i}{\bf1}_{\{\Xi_i\in A\}}\right].
        \]
        若其方差为零, 则其中心化分布集中在零点. 由多对一公式,
        \[
            \mathbb E\left[\sum_i e^{-\Xi_i}{\bf1}_{\{\Xi_i\ne0\}}\right]=0,
        \]
        因此几乎必然所有子代都位于零点. 于是
        \[
            \mathbb E[\#\Xi]=\mathbb E\left[\sum_i e^{-\Xi_i}\right]=1,
        \]
        这与 $\psi(0)>0$ 矛盾.

        固定 $a\in[0,1]$, $t>0$ 以及目标值
        $\zeta<\mu$, 其中 $\mu=-a+\frac{\pi^2}2\sigma^2t>-a$.
        当 $b=1$ 且 $d=\frac{\pi^2}{2}\sigma^2$ 时,
        $a+\zeta-dtb^{-2}=\zeta-\mu<0$. 由连续性, 可取充分接近 $1$
        且大于 $1$ 的 $b$, 以及充分接近
        $\frac{\pi^2}{2}\sigma^2$ 且小于它的 $d$, 使得
        $a+\zeta-dtb^{-2}<0$ 仍成立.
    }
    \begin{align*}
        \mathbb P\left(
            \mathcal M^{N|a}_{\ell_N}
            \le
            \zeta\log N
        \right)
        &
        =
        \mathbb P\left(
            \mathcal M^{N|a}_{\ell_N}
            \le
            \zeta\log N,\,
            F_N
        \right)
        +
        \mathbb P\left(
            \mathcal M^{N|a}_{\ell_N}
            \le
            \zeta\log N,\,
            F_N^{\mathrm c}
        \right)
        \\
        &
        \le
        \mathbb P\left(
            F_N
        \right)
        +
        \mathbb P\left(
            \mathcal M^{N|a}_{\ell_N}
            \le
            \zeta\log N,\,
            F_N^{\mathrm c}
        \right),
    \end{align*}
    \bilingual{
        and we proceed to calculate, by union bound and \nameref{lemma:many-to-one} for $\mathcal V$. We keep the endpoint inequality non-strict, since no atomlessness assumption is made, and retain $i=0$ in the drawdown minimum as required by \Cref{lemma:2.3}.
    }{
        于是通过并集界和 $\mathcal V$ 的\nameref{lemma:many-to-one}计算第二项的上界. 由于没有假设增量律无原子, 端点不等式保持非严格; 同时, 回撤最小值保留 $i=0$, 以符合\Cref{lemma:2.3}的条件.
    }
    \begin{align*}
        &\quad
        \mathbb P\left(
            \mathcal M^{N|a}_{\ell_N}
            \le
            \zeta\log N,\,
            F_N^{\mathrm c}
        \right)
        \\
        &
        =
        \mathbb E\!\left[
            \begin{aligned}
                &{\bf 1}\left\{
                    \exists k\in[\![1,\lfloor N^a\rfloor]\!],\,
                    \exists|u|=\ell_N,\,
                    \mathcal V^{N|a|(k)}(u)\le\zeta\log N
                \right\}\\
                &{\bf 1}\left\{
                    \forall j,n\in[\![0,\ell_N]\!],\,
                    \forall|u|=n,\,
                    j\le n,\,
                    \mathcal V^{N|a}(u_j)-\mathcal V^{N|a}(u)
                    \le
                    b\log N
                \right\}
            \end{aligned}
        \right]
        \\ &
        =
        \mathbb E\!\left[
            \begin{aligned}
                &{\bf 1}\left\{
                    \begin{aligned}
                        &\exists k\in[\![1,\lfloor N^a\rfloor]\!],\,
                        \exists|u|=\ell_N,\,
                        \mathcal V^{N|a|(k)}(u)\le\zeta\log N,
                        \\
                        &\forall i,j\in[\![0,\ell_N]\!],\,
                        i\le j\,\Rightarrow\,
                        \mathcal V^{N|a|(k)}(u_j)
                        -\mathcal V^{N|a|(k)}(u_i)
                        \ge
                        -b\log N
                    \end{aligned}
                \right\}\\
                &{\bf 1}\left\{
                    \forall j,n\in[\![0,\ell_N]\!],\,
                    \forall|u|=n,\,
                    j\le n\,\Rightarrow\,
                    \mathcal V^{N|a}(u_j)-\mathcal V^{N|a}(u)
                    \le
                    b\log N
                \right\}
            \end{aligned}
        \right]
        \\ &
        =
        \mathbb E\!\left[
            \begin{aligned}
                &{\bf 1}\left\{
                    \begin{aligned}
                        &\exists k\in[\![1,\lfloor N^a\rfloor]\!],\,
                        \exists|u|=\ell_N,\,
                        \mathcal V^{N|a|(k)}(u)\le\zeta\log N,
                        \\
                        &\min_{0\le i\le j\le\ell_N}\left(
                            \mathcal V^{N|a|(k)}(u_j)
                            -\mathcal V^{N|a|(k)}(u_i)
                        \right)
                        \ge
                        -b\log N
                    \end{aligned}
                \right\}\\
                &{\bf 1}\left\{
                    \forall j,n\in[\![0,\ell_N]\!],\,
                    \forall|u|=n,\,
                    j\le n\,\Rightarrow\,
                    \mathcal V^{N|a}(u_j)-\mathcal V^{N|a}(u)
                    \le
                    b\log N
                \right\}
            \end{aligned}
        \right]
        \\ &
        \le
        \sum_{k=1}^{\lfloor N^a\rfloor}
        \mathbb E\!\left[
            \sum_{|u|=\ell_N}
            {\bf 1}\left\{
                \begin{aligned}
                    \mathcal V^{N|a|(k)}(u)\le\zeta\log N,
                    \\
                    \min_{0\le i\le j\le\ell_N}\left(
                        \mathcal V^{N|a|(k)}(u_j)
                        -\mathcal V^{N|a|(k)}(u_i)
                    \right)
                    \ge
                    -b\log N
                \end{aligned}
            \right\}
        \right]
        \\ &
        =
        \sum_{k=1}^{\lfloor N^a\rfloor}
        \mathbb E\!\left[
            e^{Y_{\ell_N}}
            {\bf 1}\left\{
                \begin{aligned}
                    Y_{\ell_N}\le\zeta\log N,
                    \\
                    \min_{0\le i\le j\le\ell_N}\left(
                        Y_j
                        -Y_i
                    \right)
                    \ge
                    -b\log N
                \end{aligned}
            \right\}
        \right]
        \\ &
        \le
        N^{a+\zeta}
        \mathbb P\left(
            Y_{\ell_N}\le\zeta\log N,\,
            \min_{0\le i\le j\le\ell_N}\left(
                Y_j
                -Y_i
            \right)
            \ge
            -b\log N
        \right)
        \\ &
        \le
        N^{a+\zeta}
        \mathbb P\left(
            Y_{\ell_N}<(b+\zeta+a)\log N,\,
            \min_{0\le i\le j\le\ell_N}\left(
                Y_j
                -Y_i
            \right)
            \ge
            -b\log N
        \right)
        \\ &
        \le
        N^{a+\zeta}
        e^{-d\frac{\ell_N}{(b\log N)^2}}
        =
        N^{a+\zeta-dt b^{-2}+o_N(1)},
    \end{align*}
    \bilingual{
        Choose $c\in(0,\frac12)$ small enough that $b+a+\zeta<b/c$; this is
        possible because $b/c\to\infty$ as $c\downarrow0$. The endpoint bound
        in the preceding display is then at most $b\log N/c$, so the endpoint
        event is contained in the event of \Cref{lemma:2.3}.
        Since $\ell_N/\log^3N\to t$ and $2+c<3$, we have
        $\ell_N\ge(b\log N)^{2+c}$ for all sufficiently large $N$. Thus the
        last inequality follows from \Cref{lemma:2.3} with
        $\Delta=b\log N$ and $L=\ell_N$. The choices of $b,d$ above give
        $a+\zeta-dtb^{-2}<0$. Hence,
    }{
        取充分小的 $c\in(0,\frac12)$ 使得 $b+a+\zeta<b/c$; 这是因为
        $c\downarrow0$ 时 $b/c\to\infty$. 于是上一个公式中的端点上界不超过
        $b\log N/c$, 故端点事件包含于\Cref{lemma:2.3}中的事件. 由于
        $\ell_N/\log^3N\to t$ 且 $2+c<3$, 对充分大的 $N$ 有
        $\ell_N\ge(b\log N)^{2+c}$. 因此, 对参数
        $\Delta=b\log N$ 和 $L=\ell_N$, 由\Cref{lemma:2.3} 得到最后一个
        不等式. 上文选择的 $b,d$ 满足 $a+\zeta-dtb^{-2}<0$. 因此,
    }
    \begin{equation*}
        \mathbb P\left(
            \mathcal M^{N|a}_{\ell_N}
            \le
            \zeta\log N
        \right)
        \le
        N^{-(b-1)+o_N(1)}
        +
        N^{a+\zeta-dt b^{-2}+o_N(1)}
        \le
        N^{-\varepsilon},
    \end{equation*}
    \bilingual{
        where we choose some $\varepsilon\in(0,\frac 12\min\{b-1,dtb^{-2}-a-\zeta\})$ such that $\varepsilon<b-1$ and $\varepsilon<dtb^{-2}-a-\zeta$, and it holds for some $N_0=N_0(a,t,\zeta)>N_1$ and all $N\ge N_0$. This proves \eqref{eq:4.16} for all $a\in[0,1]$.
    }{
        其中, 选择某个 $\varepsilon\in(0,\frac 12\min\{b-1,dtb^{-2}-a-\zeta\})$ 以使得 $\varepsilon<b-1$ 且 $\varepsilon<dtb^{-2}-a-\zeta$, 则对于某个 $N_0=N_0(a,t,\zeta)>N_1$ 和所有 $N\ge N_0$ 成立. 至此, 对所有 $a\in[0,1]$ 的情形, \eqref{eq:4.16} 得证.
    }
\end{proof}

\begin{lemma}[\bilingual{One-sided $L^1$ endpoint estimate at $a=1$}{\(a=1\) 时端点的一侧 $L^1$ 估计}]
    \label{lemma:left-tail-L1-endpoint}
    \bilingual{
        Assume the hypotheses listed in \Cref{theorem:speed-L1}, with no
        fourth-moment assumption. Fix $t>0$ and an integer generation sequence
        $\ell_N$ such that
        $\ell_N/\log^3 N\to t$. Then
    }{
        假设\Cref{theorem:speed-L1} 中列出的条件, 不要求四阶矩. 固定
        $t>0$ 及满足 $\ell_N/\log^3 N\to t$ 的整数代数列, 则
    }
    \begin{equation*}
        \liminf_{N\to\infty}
        \frac{\mathbb E[\mathcal M^{N|1}_{\ell_N}]}{\log N}
        \ge -1+\frac{\pi^2}{2}\sigma^2t.
    \end{equation*}
\end{lemma}

\begin{proof}
    \bilingual{
        Fix any $\zeta<-1+\frac{\pi^2}{2}\sigma^2t$. The lower-tail estimate
        \eqref{eq:4.16}, whose proof above uses no fourth moment, gives some
        $\varepsilon>0$ such that for all sufficiently large $N$,
    }{
        任取 $\zeta<-1+\frac{\pi^2}{2}\sigma^2t$. 上文证明的左尾估计
        \eqref{eq:4.16} 不使用四阶矩, 因而存在 $\varepsilon>0$, 使得对
        所有充分大的 $N$,
    }
    \[
        p_N:=\mathbb P(\mathcal M^{N|1}_{\ell_N}\le\zeta\log N)
        \le N^{-\varepsilon}.
    \]
    \bilingual{
        The cited moment lemmas imply that $\mathcal M^{N|1}_{\ell_N}$ is
        integrable. Also, \Cref{lemma:M:exp-neg-moment} gives
        $\mathbb E[e^{-\mathcal M^{N|1}_{\ell_N}}]\le N$. Set
        $L_N=\log N$ and
        $A_N=\{\mathcal M^{N|1}_{\ell_N}\le\zeta L_N\}$. Since
        $x\ge-e^{-x}$, on $A_N$ we have
        \[
            \mathcal M^{N|1}_{\ell_N}
            \ge -L_N-e^{-L_N}e^{-\mathcal M^{N|1}_{\ell_N}}.
        \]
        On $A_N^c$ the endpoint exceeds $\zeta L_N$. Splitting its
        expectation gives
    }{
        所引矩引理保证 $\mathcal M^{N|1}_{\ell_N}$ 可积. 此外,
        \Cref{lemma:M:exp-neg-moment} 给出
        $\mathbb E[e^{-\mathcal M^{N|1}_{\ell_N}}]\le N$. 令
        $L_N=\log N$ 且 $A_N=\{\mathcal M^{N|1}_{\ell_N}\le\zeta L_N\}$.
        由 $x\ge-e^{-x}$, 在 $A_N$ 上有
        \[
            \mathcal M^{N|1}_{\ell_N}
            \ge -L_N-e^{-L_N}e^{-\mathcal M^{N|1}_{\ell_N}}.
        \]
        在 $A_N^c$ 上端点大于 $\zeta L_N$. 按这两个事件分解期望, 得
    }
    \begin{equation*}
        \mathbb E[\mathcal M^{N|1}_{\ell_N}]
        \ge
        \zeta L_N
        -\bigl(\max\{\zeta,0\}+1\bigr)L_Np_N
        -1.
    \end{equation*}
    \bilingual{
        The positive-part term is needed when $\zeta<0$: in that case the
        good-event contribution is bounded below by $\zeta L_N$, not by
        $\zeta L_N(1-p_N)$. Divide by $L_N$ and let $N\to\infty$ to obtain
        $\liminf_N L_N^{-1}\mathbb E[\mathcal M^{N|1}_{\ell_N}]\ge\zeta$.
        Since this holds for every $\zeta$ below the target, the claimed
        endpoint estimate follows. This argument uses only the lower-tail
        probability and the exponential bound on the negative part; it does
        not use the $L^2$ trajectory conclusion of Proposition 3.1.
    }{
        当 $\zeta<0$ 时必须保留正部项: 此时好事件上的期望下界为
        $\zeta L_N$, 而非 $\zeta L_N(1-p_N)$. 两边除以 $L_N$, 令
        $N\to\infty$, 得到
        $\liminf_N L_N^{-1}\mathbb E[\mathcal M^{N|1}_{\ell_N}]\ge\zeta$.
        由于这对所有低于目标值的 $\zeta$ 都成立, 所述端点估计得证.
        该论证只使用左尾概率估计和负部的指数矩界, 不使用命题 3.1
        的 $L^2$ 轨迹结论.
    }
\end{proof}
